% ==============================================================================
% Final Journal Figure Captions for Figure 1 and Figure 2
% Target: Software: Practice and Experience (Wiley)
% ==============================================================================

% ------------------------------------------------------------------------------
% Figure 1 Caption
% ------------------------------------------------------------------------------
\begin{figure*}[!t]
\centering
\includegraphics[width=\textwidth]{figs/figure1_secchaos_framework.pdf}
\caption{Overall SecChaos-RL closed-loop resilience engineering framework across five coordinated functional regions: 
(Region~A) Monitored microservice dependency DAG $\mathcal{G} = (\mathcal{V}, \mathcal{E})$ with runtime telemetry probes; 
(Region~B) Security Chaos Engine injecting standardized stochastic faults ($F_1$--$F_8$, Poisson arrival rate $\lambda_k \in [0.05, 0.25]$) alongside dynamic continuous logistic threat severity evaluation ($\mathcal{L}_{\mathrm{threat}}$); 
(Region~C) Threat-aware deep reinforcement learning controller featuring linear state projection, threat-aware feature gating, 2-head self-attention ($H=2$), and three decoupled output heads (actor policy $\pi_\theta$ with adaptive temperature $\tau_t$, critic baseline $V_\phi$, and auxiliary risk prediction head $\hat{\mathbf{r}}_{t+1}$); 
(Region~D) Operational safety guardrails (dispatch validation and transient queue drop filter $\epsilon_{\mathrm{queue}} = 0.08$), asynchronous execution queue ($\tau \in \{0, 1, 2\}$), and cloud-native remediation actuators delivering active service healing (reducing targeted service risk by $80\%\sigma$ and restoring availability by $+15\%\sigma$); and 
(Region~E) Real-time multi-objective evaluation tracking availability, recovery time (MTTR), security risk reduction, and false mitigation rate (FMR), driving PPO policy gradient parameter updates ($\nabla_\theta \mathcal{L}_{\mathrm{PPO}}$) through the closed learning loop. 
Solid lines indicate runtime data and operational control paths; dashed red lines indicate chaos fault ingress; and dotted purple lines indicate policy-training updates (performed during offline training phases). 
Abbreviations: DAG: Directed Acyclic Graph; DRL: Deep Reinforcement Learning; FMR: False Mitigation Rate; MTTR: Mean Time to Recovery; PPO: Proximal Policy Optimization; SCE: Security Chaos Engineering.}
\label{fig:framework_architecture}
\end{figure*}

% ------------------------------------------------------------------------------
% Figure 2 Caption
% ------------------------------------------------------------------------------
\begin{figure*}[!t]
\centering
\includegraphics[width=0.96\textwidth]{figs/figure2_secchaos_policy.pdf}
\caption{Threat-aware deep neural policy network architecture and optimization workflow: 
(Stage~A) Standardized 16-dimensional observation state $\mathbf{s}_t \in \mathbb{R}^{16}$ partitioned into five semantic feature groups (security risks $s_0$--$s_4$, threat dynamics $s_5$--$s_7$, cluster resource telemetry $s_8$--$s_{10}$, resilience QoS $s_{11}$--$s_{12}$, and execution status $s_{13}$--$s_{15}$); 
(Stage~B) Sixteen independent linear feature projections transforming scalar inputs into $32$-dimensional token embeddings ($\mathbf{e}_t^{(i)} \in \mathbb{R}^{32}$); 
(Stage~C) Context-aware threat gating multiplying security channels by $(1.0 + \mathcal{L}_{\mathrm{threat}})$ based on global threat severity $s_5$; 
(Stage~D) Multi-head self-attention ($H=2, d_k=16$) capturing topological inter-service failure dependencies across tokens; 
(Stage~E) Flattening and shared dense representation layers ($512 \to 256 \to 128$, Tanh); 
(Stage~F) Three multi-task decoupled output heads: Actor policy $\pi_\theta(a|\mathbf{s}_t)$ with threat-adaptive exploration temperature $\tau_t = \tau_0 \exp(-\beta \mathcal{L}_{\mathrm{threat}})$, Critic baseline $V_\phi(\mathbf{s}_t)$, and Auxiliary next-step risk predictor $\hat{\mathbf{r}}_{t+1} \in \mathbb{R}^6$; and 
(Stage~G) PPO policy optimization with trajectory rollouts, clipped surrogate objective, value loss, and auxiliary risk prediction loss. 
The auxiliary head acts as an inductive regularizer during training and is detached during live inference ($0.44$\,ms online decision latency). 
Abbreviations: GAE: Generalized Advantage Estimation; MHA: Multi-Head Attention; MLP: Multi-Layer Perceptron; QoS: Quality of Service.}
\label{fig:network_architecture}
\end{figure*}
